\documentstyle[%


righttag%


,11pt%


,amssymb%


,russian%


,izvru%


]{amsart}





\newcommand{\tl}[3]{\medskip\noindent{\bf#1}\ #2 \dotfill #3 \par} 


\pagestyle{empty}


% \pagestyle{plain}


\begin{document}


% \setcounter{page}{80}


\par


\vskip 2cm


\par


\bigskip


\bigskip


\bigskip


\par


\centerline{\Large СОДЕРЖАНИЕ}


\bigskip


\bigskip


%%%%%%%%%% Use \newline %%%%%%%%%%%





\tl{Булгаков А.И., Ткач Л.И.}


{Возмущение однозначного оператора многозначным\newline отображением


типа Гаммерштейна с невыпуклыми образами}{3}


\tl{Вагабов А.И.} 


{Задача о колебании конечной струны 


с нелинейным возмущением}{17}


\tl{Джохадзе О.М.}


{Задача типа Дарбу в трехгранном угле для


уравнения третьего порядка\newline гиперболического типа}{22}


\tl{Исламов Г.Г.}


{О полиэдральной разрешимости системы 


линейных дифференциальных\newline уравнений}{31}


\tl{Коняев Ю.А.} 


{Итерационный метод анализа нелинейных 


сингулярно возмущенных\newline начальных 


и краевых задач}{38}


\tl{Нахман А.Д.} 


{О частных суммах кратных рядов Фурье 


по многочленам Якоби}{46}


\tl{Широков И.В.} 


{Исследование устойчивости решений дифференциальных


уравнений,\newline допускающих транзитивную группу симметрии}{57}


\tl{Якупов М.Ш.}  


{Каноническое описание и обращение операторов 


в дискретном\newline пространстве. II}{64}





\bigskip


\par


\centerline{КРАТКИЕ СООБЩЕНИЯ}


\bigskip





\tl{Годуля Я., Старков В.В.}


{К теореме регулярности 


для линейно-инвариантных семейств\newline функций в поликруге}{75}


\tl{Михайлов А.Б.} 


{О двухточечной задаче для уравнения второго порядка 


в частных\newline производных с постоянными коэффициентами}{80}








\vfill


\noindent\raisebox{2pt}{\copyright} Казанский госудаpственный унивеpситет


\end{document}


